\documentclass[11pt]{article}
\usepackage[a4paper,margin=1in]{geometry}
\usepackage[T1]{fontenc}
\usepackage[utf8]{inputenc}
\usepackage{lmodern,microtype}
\usepackage{amsmath,amssymb,amsthm,mathtools}
\usepackage{enumitem,xcolor,hyperref}
\usepackage[nameinlink,capitalise,noabbrev]{cleveref}
\hypersetup{colorlinks=true,linkcolor=blue!60!black,citecolor=blue!60!black,urlcolor=blue!60!black}

\newtheorem{definition}{Definition}[section]
\newtheorem{theorem}[definition]{Theorem}
\newtheorem{proposition}[definition]{Proposition}
\newtheorem{corollary}[definition]{Corollary}
\newtheorem{conjecture}[definition]{Conjecture}
\newtheorem{remark}[definition]{Remark}

\newcommand{\C}{\mathbb C}
\newcommand{\Oterm}{\mathcal O}
\newcommand{\D}{\mathrm D}
\newcommand{\Q}{\mathcal Q}
\newcommand{\norm}[1]{\left\lVert #1\right\rVert}
\DeclareMathOperator{\Sym}{Sym}

\title{All-Order Tensor Inverse Jets for Multidimensional Pandrosion\\[3pt]
\large Multilinear normal forms, $m$-flat charts, and anchor-order multiplier suppression}
\author{Ivan Besevic\\Independent Mathematical Research}
\date{May 2026}

\begin{document}
\maketitle

\begin{abstract}
The scalar higher-order Pandrosion theory shows that local inverse jets flatten an analytic equation and make the fixed-anchor multiplier scale like a high power of the anchor distance. The tensorial multidimensional paper proved the first nontrivial case: a quadratic inverse jet cancels the bilinear residual tensor of an analytic system. This paper gives the all-order version. Let \(G:\C^n\to\C^n\) be analytic near a regular zero \(\rho\), with \(A=\D G(\rho)\) invertible. For every \(m\ge1\), there is a unique polynomial jet
\[
        \phi_m(u)=\rho+B_1(u)+B_2(u,\ldots,u)+\cdots+B_m(u,\ldots,u)
\]
whose homogeneous symmetric multilinear coefficients satisfy
\[
        G(\phi_m(u))=u+\Oterm(\norm{u}^{m+1}).
\]
The coefficients are determined recursively: at each order only \(A B_k\) is unknown, and all lower terms are already fixed. Thus \(\phi_m\) is the degree-\(m\) Taylor jet of the local inverse of \(G\). For the charted averaged-Jacobian Pandrosion map, this gives
\[
        \norm{D P_{G\circ\phi_m,a}(0)}=\Oterm(\norm{a}^{m}).
\]
More sharply, if the first uncancelled defect is a homogeneous tensor \(H_{m+1}\), then
\[
        D P_{G\circ\phi_m,a}(0)
        =
        \frac{1}{m+1}D H_{m+1}(a)+\Oterm(\norm{a}^{m+1}).
\]
The paper closes with a finite-probe reconstruction conjecture: a stable Pandrosion engine that estimates multilinear defects from symmetric probes should recover the inverse-jet hierarchy asymptotically.
\end{abstract}

\paragraph{Scope and status.}
This is a theorem-and-conjecture manuscript inside the Pandrosion research programme. The theorems are local analytic normal-form statements; the final reconstruction statement is conjectural and meant to be tested against finite-slope engines.

\section{From tensor-flat charts to all orders}

Earlier scalar papers established that the anchored divided-difference map improves when the equation is first expressed in a chart where the transformed residual is nearly linear. The multidimensional tensor paper identified the leading obstruction:
\[
        F(u)=u+\frac12 T_2[u,u]+\cdots
        \quad\Longrightarrow\quad
        P_{F,a}(s)=\frac12 T_2[a,s]+\cdots .
\]
Cancelling \(T_2\) by a quadratic inverse jet makes the multiplier second order in the anchor distance. The natural continuation is to cancel every homogeneous tensor up to a prescribed order.

\section{Homogeneous tensor jets}

Let \(G\) be analytic near \(\rho\in\C^n\), \(G(\rho)=0\), and \(A=\D G(\rho)\) invertible. Write the Taylor expansion
\begin{equation}\label{eq:G-expansion}
        G(\rho+w)
        =
        A w+\sum_{q\ge2} G_q[w,\ldots,w],
\end{equation}
where \(G_q\in \Sym^q((\C^n)^*)\otimes\C^n\) is the symmetric \(q\)-linear coefficient \(G_q=\D^qG(\rho)/q!\).

\begin{definition}[Polynomial inverse jet]
An order-\(m\) tensor jet at \(\rho\) is a polynomial map
\[
        \phi_m(u)=\rho+\sum_{k=1}^m B_k[u,\ldots,u],
\]
where \(B_k\in \Sym^k((\C^n)^*)\otimes\C^n\). It is \(m\)-flat for \(G\) if
\[
        G(\phi_m(u))=u+\Oterm(\norm{u}^{m+1}).
\]
\end{definition}

\begin{theorem}[All-order tensor inverse jet]\label{thm:inverse-jet}
For every \(m\ge1\), there is a unique order-\(m\) tensor jet \(\phi_m\) such that
\[
        G(\phi_m(u))=u+\Oterm(\norm{u}^{m+1}).
\]
The coefficients are determined recursively. First \(B_1=A^{-1}\). If \(B_1,\ldots,B_{k-1}\) are known, let \(R_k(u)\) be the homogeneous degree-\(k\) part of
\[
        \sum_{q=2}^{k}G_q[W_{k-1}(u),\ldots,W_{k-1}(u)],
        \qquad
        W_{k-1}(u)=\sum_{j=1}^{k-1}B_j[u,\ldots,u].
\]
Then
\[
        B_k[u,\ldots,u]=-A^{-1}R_k(u).
\]
Moreover, \(\phi_m\) is exactly the degree-\(m\) Taylor jet of the local inverse of \(G\) at \(0\).
\end{theorem}

\begin{proof}
The inverse function theorem gives a local inverse \(\Psi\) with \(G(\Psi(u))=u\), so existence of a formal Taylor jet is guaranteed. The recursive proof also gives uniqueness. For \(k=1\), the degree-one term in \(G(\phi_m(u))\) is \(AB_1u\), hence \(B_1=A^{-1}\).

Assume \(B_1,\ldots,B_{k-1}\) have been chosen so that all homogeneous terms of degrees \(2,\ldots,k-1\) vanish. In degree \(k\), the only term involving the unknown \(B_k\) is \(AB_k[u,\ldots,u]\), because every nonlinear \(G_q\) has \(q\ge2\) and therefore uses only the already known lower coefficients to contribute to degree \(k\). The remaining degree-\(k\) contribution is \(R_k(u)\). The equation
\[
        AB_k[u,\ldots,u]+R_k(u)=0
\]
has the unique solution stated above because \(A\) is invertible. Induction constructs a unique jet through order \(m\), and it must agree with the Taylor truncation of the analytic inverse.
\end{proof}

\begin{remark}
For \(m=2\), \cref{thm:inverse-jet} gives
\[
        \phi_2(u)=\rho+A^{-1}u-A^{-1}G_2[A^{-1}u,A^{-1}u],
\]
which is the quadratic tensor-flat chart from the multidimensional Pandrosion paper.
\end{remark}

\section{Averaged-Jacobian Pandrosion in an \(m\)-flat chart}

Let \(F:\C^n\to\C^n\) be analytic near \(0\), \(F(0)=0\), and \(DF(0)=I\). The averaged-Jacobian divided difference is
\[
        \Q_F(a,s)=\int_0^1 DF(s+t(a-s))\,dt,
\]
so that \(F(a)-F(s)=\Q_F(a,s)(a-s)\). When \(\Q_F(a,s)\) is invertible, the anchored multidimensional Pandrosion map is
\[
        P_{F,a}(s)=a-\Q_F(a,s)^{-1}F(a)
        =
        s-\Q_F(a,s)^{-1}F(s).
\]

\begin{definition}[\(m\)-flat residual]
The residual \(F\) is \(m\)-flat at \(0\) if
\[
        F(u)=u+\Oterm(\norm{u}^{m+1}).
\]
If
\[
        F(u)=u+H_{m+1}(u)+\Oterm(\norm{u}^{m+2}),
\]
with \(H_{m+1}\) homogeneous of degree \(m+1\), then \(H_{m+1}\) is the leading defect tensor.
\end{definition}

\begin{theorem}[Anchor-order multiplier law]\label{thm:mult-law}
Assume \(F\) is \(m\)-flat at \(0\). Then, for small anchors \(a\) with \(\Q_F(a,0)\) invertible,
\[
        \norm{D P_{F,a}(0)}=\Oterm(\norm{a}^{m}).
\]
If the leading defect is \(H_{m+1}\), then
\begin{equation}\label{eq:leading-mult}
        D P_{F,a}(0)
        =
        \frac{1}{m+1}D H_{m+1}(a)+\Oterm(\norm{a}^{m+1}).
\end{equation}
\end{theorem}

\begin{proof}
Using the secant form of the map and \(F(0)=0\),
\[
        D P_{F,a}(0)=I-\Q_F(a,0)^{-1}DF(0)=I-\Q_F(a,0)^{-1}.
\]
Since \(F(u)=u+\Oterm(\norm{u}^{m+1})\), we have \(DF(u)=I+\Oterm(\norm{u}^{m})\), hence
\[
        \Q_F(a,0)=\int_0^1 DF(ta)\,dt=I+\Oterm(\norm{a}^{m}).
\]
This proves the norm estimate. If \(F(u)=u+H_{m+1}(u)+\Oterm(\norm{u}^{m+2})\), then
\[
        DF(ta)=I+t^m DH_{m+1}(a)+\Oterm(\norm{a}^{m+1}),
\]
so
\[
        \Q_F(a,0)=I+\frac{1}{m+1}DH_{m+1}(a)+\Oterm(\norm{a}^{m+1}).
\]
Expanding the inverse gives \eqref{eq:leading-mult}.
\end{proof}

\begin{corollary}[All-order charted Pandrosion]\label{cor:charted}
Let \(G(\rho)=0\) and \(DG(\rho)\) be invertible. Let \(\phi_m\) be the all-order tensor inverse jet from \cref{thm:inverse-jet}, and set \(F_m=G\circ\phi_m\). Then
\[
        \norm{D P_{F_m,a}(0)}=\Oterm(\norm{a}^{m}).
\]
\end{corollary}

\begin{proof}
By construction, \(F_m(u)=u+\Oterm(\norm{u}^{m+1})\). Apply \cref{thm:mult-law}.
\end{proof}

\section{Two-point residual estimate}

\begin{proposition}[Uniform two-point order]\label{prop:two-point}
If \(F(u)=u+\Oterm(\norm{u}^{m+1})\), then for \(a,s\) sufficiently small,
\[
        \norm{P_{F,a}(s)}
        \le
        C\norm{s}\,(\norm{a}+\norm{s})^m.
\]
\end{proposition}

\begin{proof}
Write \(F(s)=s+R(s)\), with \(\norm{R(s)}\le C_0\norm{s}^{m+1}\). Also
\[
        \Q_F(a,s)=I+E(a,s),
        \qquad
        \norm{E(a,s)}\le C_1(\norm{a}+\norm{s})^m.
\]
For \(a,s\) small, \(\Q_F(a,s)^{-1}=I-E(a,s)+\Oterm(\norm{E}^2)\). Therefore
\[
        P_{F,a}(s)=s-\Q_F(a,s)^{-1}F(s)
        =
        E(a,s)s-R(s)+\Oterm(\norm{s}\norm{E(a,s)}^2),
\]
which gives the claimed bound after shrinking the neighborhood.
\end{proof}

\section{Finite-probe reconstruction conjecture}

The theorems above use derivatives to define the inverse jet. The practical Pandrosion programme uses finite probes. The bridge is a reconstruction conjecture.

\begin{conjecture}[Finite-probe inverse-jet reconstruction]
Let \(G\) be analytic near a regular zero. Suppose a Pandrosion engine uses symmetric probe clouds whose radius tends to zero, rejects rank-defective finite-slope matrices, and estimates homogeneous defects by multilinear polarization identities. Then the estimated chart coefficients converge order by order to the tensor inverse-jet coefficients \(B_1,\ldots,B_m\) of \cref{thm:inverse-jet}.
\end{conjecture}

\begin{remark}
This conjecture is deliberately algorithmic. The theorem says what the optimal tensors are; the conjecture says that a finite-slope engine can recover them without direct derivative access.
\end{remark}

\section{Conclusion}

The all-order tensor inverse jet is the multidimensional normal form behind Pandrosion chart optimization. It unifies scalar higher-order cancellation, quadratic tensor-flat charts, and the atlas programme: every regular zero has a canonical hierarchy of local coordinates in which the raw anchored multiplier is suppressed from \(\Oterm(\norm{a})\) to \(\Oterm(\norm{a}^m)\). The next step is to make these local charts uniform across a finite atlas.

\begin{thebibliography}{9}
\bibitem{general}
I. Besevic, \emph{Anchored Pandrosion for General Analytic Equations}, manuscript, May 2026.
\bibitem{jets}
I. Besevic, \emph{Higher-Order Jet Cancellation for Anchored Divided-Difference Iterations}, manuscript, May 2026.
\bibitem{tensor}
I. Besevic, \emph{Tensorial Multidimensional Pandrosion}, manuscript, May 2026.
\bibitem{atlas}
I. Besevic, \emph{Stable Numerical Atlases for Pandrosion Dynamics}, manuscript, May 2026.
\end{thebibliography}

\end{document}
